The implementation retains integrated CP rooming and greedy post-processing as baselines, and adds \texttt{room\_mode=decomposed} as the research method.
Strict CP rooming now uses every eligible room by default; candidate truncation is explicit and opt-in.
For a co-located cluster, capacity is computed over the union of participating groups rather than per activity.

\subsection{Fixed-Time Room Jobs}
Given a master schedule, co-located activities are collapsed into room jobs $j\in\mathcal{J}$.
Each job has a fixed week, day, start, duration, union robust-demand requirement $D_j$, and candidate domain $R_j$.
The domain is the intersection of member eligibility after room type, specialization, capacity (when hard), availability, closures, and room locks are applied.
Binary $y_{j,r}$ selects exactly one candidate room:
\begin{align}
\sum_{r\in R_j}y_{j,r}=1 \qquad \forall j\in\mathcal{J}.
\end{align}
Per-room/per-slot at-most-one constraints enforce physical occupancy.
Additional pairwise constraints enforce travel buffers for jobs sharing staff or groups, and equality constraints implement an optional repeated weekly room pattern.
The room objective prioritizes tutorial room type, capacity overflow when capacity is soft, capacity waste, and deterministic room identifiers.
In objective-free runs the exact subproblem stops at feasibility rather than spending the remaining budget proving room-waste optimality; objective-enabled runs retain the room objective.

paragraph{Support-aware feedback proxy.}
For a fixed incumbent-majority room $r_c$ for each course $c$, the compatibility feedback scalar counts event-level collisions at $(q,r_c)$.
The opt-in ablation instead resolves every period--majority-room collision deterministically, retains one capacity-prioritized event, and counts each displaced course once over the timetable.
This distinct-course fragmentation count is only a candidate-ordering proxy: it is neither the official room-stability objective nor an exact fixed-time room solve.
All resulting candidates are lifted to physical rooms and accepted only after full hard validation and canonical official rescoring.

\subsection{Hall-Deficiency Certificates}
For occupied slot $q=(w,d,s)$, form a bipartite graph between active jobs $\mathcal{J}_q$ and their eligible rooms.
A maximum matching is computed.
If it does not cover every active job, alternating reachability from unmatched jobs yields a witness $J'\subseteq\mathcal{J}_q$ with neighborhood $N(J')$ satisfying
\begin{align}
|J'|>|N(J')|.
\end{align}
By Hall's theorem, those jobs cannot receive distinct rooms at $q$.
The certificate stores the activities, representative cluster activities, candidate rooms, slot, deficiency, active demand mode, required capacities, and a counting explanation.

\subsection{Master Cut}
Let $D(j,t)$ be the exact candidate-room domain of representative job $j$ when started at $t$, after cluster membership, aggregate demand, room type, specialization, room locks, interval availability, and closures are applied.
Retain the fixed-time certificate's witnessed room set $R^{\mathrm w}$ as evidence, and define the option-expanded literal set
\begin{align}
U(q,R^{\mathrm w})=
\{(j,t):q\in I(j,t),\;D(j,t)\subseteq R^{\mathrm w}\}.
\end{align}
Only starts whose complete candidate domain is contained in the witnessed room set are counted.
Their actual neighborhood is reconstructed separately as
\begin{align}
\Gamma(U)=\bigcup_{(j,t)\in U(q,R^{\mathrm w})}D(j,t).
\end{align}
The master receives the standard option-expanded Hall inequality
\begin{align}
\sum_{(j,t)\in U(q,R^{\mathrm w})}x_{j,t}
\le |\Gamma(U)|. \label{eq:contextual-hall}
\end{align}
The witness $R^{\mathrm w}$ therefore limits which alternatives can enter the derivation, while the right-hand side is the stronger reconstructed neighborhood cardinality, not $|R^{\mathrm w}|$.
For any selected subset of those alternatives, distinct active jobs require distinct rooms drawn from a subset of $\Gamma(U)$, which proves~\eqref{eq:contextual-hall}; the master's at-most-one start constraint prevents two alternatives of one job from being selected together.
The incumbent witnessed starts must all satisfy the subset rule, and the contextual inequality is emitted only when it excludes the incumbent and is stronger than its exact incumbent-start nogood.
Non-incumbent starts whose domain expands outside $R^{\mathrm w}$ are excluded from $U$.
If an incumbent domain is not contained in $R^{\mathrm w}$, if any domain proof is uncertain, or if infeasibility arises only after travel or repeated-room coupling, the solver emits the universally safe exact incumbent-start nogood instead.
Each cut records its kind, witness slot and rooms, derived $\Gamma(U)$, right-hand side, included and excluded start counts, proof rule, and fallback reason.
The master is re-solved and the exact room subproblem is repeated until a full assignment is found, the master proves infeasible, or the shared wall-clock/round limit is reached.

\subsection{Replay and Integrity Lineage}
Before adding a contextual Hall inequality, a separate checker replays the serialized certificate and derivation without consulting the constructed CP model.
It verifies the certificate schema and deficiency, fingerprints the room-domain input context, reconstructs every effective room domain from the stated members, time, duration, and master start domains, checks that the counted option set is complete, recomputes $\Gamma(U)$ and the right-hand side, and confirms that the inequality excludes the incumbent.
Malformed, incomplete, or inconsistent evidence is never inserted as a contextual cut; when the representative incumbent starts are available, the cut layer uses the exact-nogood fallback.
The replay path is independent of the cut-construction code, although both deliberately call the same authoritative effective-domain routine; it is not a second implementation of room eligibility.

Canonical JSON payloads receive content-addressed SHA-256 identifiers for the certificate, room context, inequality, and complete derivation.
The certificate-derived LNS signal retains these identifiers, the witness and derived room sets, and the domain assumptions; a selected certificate neighborhood records the source certificate, cut, and derivation identifiers in its trace.
These hashes provide deterministic integrity and lineage, not authorship or authenticity: they are neither digital signatures nor protection against an actor replacing an artifact and all of its reference hashes.

\subsection{Explanation and Repair Reuse}
The Hall certificate supports the statement ``$k$ simultaneous room jobs can use only $m<k$ rooms'' and recommends moving at least one certified job outside the slot.
For repair, certified activities are expanded through a typed activity hypergraph linking shared groups, staff, courses, rooms, clusters, distribution constraints, and temporal partition boundaries.
Activities outside the selected neighborhood are fixed through CP-SAT assumptions on a reusable exact model; an infeasible assumption core is converted into another typed search signal.
A discounted-UCB controller chooses between certificate, penalty, partition-boundary, and random families at several target sizes.
Only schedules accepted by the separate repository-local post-solve validator and found to be strictly improving replace the incumbent, and every exact-repair round records build, assumption-setup, search, validation, objective, bound, and local-optimality evidence.
Cut-guided LNS, explanation-guided repair, adaptive LNS, and bandit selection have prior art; the benefit of this integrity-linked systems integration remains an experimental question pending the prescribed ablations.
